% !TEX root = template-sojo-genin-lua.tex
% !TEX recipe = LuaLaTeX (kian-lua)
% ============================================================
%  訴状（請求の原因）テンプレート  ― IT化対応版  (kian-v08 LuaLaTeX 版)
%  用途: 裁判IT化提出用（表紙なし・請求の原因のみの単独PDF）
%  ※ 表紙あり一体型は template-sojo-lua.tex を使用
%
%  コンパイル: latexmk -r latexmkrc-lua template-sojo-genin-lua.tex
%              または VS Code の「LuaLaTeX (kian-lua)」レシピ
%  プレースホルダー: 【...】を実際の内容に置き換えてください
%  異体字: \ivs{文字} で IPAmj 明朝字形に切り替え（例: \ivs{角󠄁}）
%  ※ \SecA{} 等の {} 内に % を書くとコメントになりエラーが出ます
%
%  段落番号: このファイルは独立コンパイルのため \SecA は自動で
%  第１から開始する（請求の趣旨ファイルとの番号連続なし）。
%  見出し番号: \SecZ=番号なし \SecA=第１ \SecB=１ \SecC=⑴ \SecD=ア \SecE=(ア)
%  ※ 見出し（\SecA{} 等）は本文なしで連続させると番号表示が崩れるため、
%    各見出しの直後に【ここに記載】のプレースホルダー本文を入れています
% ============================================================

\documentclass{ltjsarticle}
\usepackage{kian-lua}                % kian-lua/ を TEXINPUTS に含めてコンパイル
% \usepackage{kian-ivs,luatexja-ruby}  % IVS 異体字（\ivs{文字}）・ルビ用

%  【注意】表紙系コマンド（\Tit \Dat \Sen \Idcs \Recs）・\maketitle は
%  置かない。当事者・裁判所はシステムのフォーム入力で処理される。

\begin{document}

% ---------------------------------------------------------- %
%  請求の原因
%  \SecZ は番号なし見出し（「第１」等の番号が付かない）
%  以降の \SecA は独立ファイルのため自動で第１から開始
% ---------------------------------------------------------- %

\SecZ{請求の原因}

% ---------------------------------------------------------- %
%  第１　当事者
% ---------------------------------------------------------- %

\SecA{当事者}



% ---------------------------------------------------------- %
%  第２　請求原因事実
% ---------------------------------------------------------- %

\SecA{請求原因事実}

\SecB{}

【ここに記載】

\SecB{}

【ここに記載】

% ---------------------------------------------------------- %
%  第３　損害
% ---------------------------------------------------------- %

\SecA{損害}

\SecB{}

【ここに記載】

\SecB{}

【ここに記載】

% ---------------------------------------------------------- %
%  第４　まとめ
% ---------------------------------------------------------- %

\SecA{まとめ}



%  【注意】証拠方法・附属書類は記載しない。
%  証拠は別途 template-shoko-lua.tex（証拠説明書）で提出する。

\vspace{2\zw}

\begin{flushright}
以　上
\end{flushright}

\end{document}
